% 图契约：本题的源、边界和推导逐项对应；箭头符号按图注解释。
\documentclass[tikz,border=7pt]{standalone}
\usepackage{fontspec,xeCJK,amsmath,amssymb}
\IfFontExistsTF{Times New Roman}{\setmainfont{Times New Roman}}{\setmainfont{TeX Gyre Termes}}
\IfFontExistsTF{Heiti SC}{\setCJKmainfont[AutoFakeBold=1.5]{Heiti SC}}{\setCJKmainfont[AutoFakeBold=1.5]{FandolHei-Regular}}
\usetikzlibrary{arrows.meta,bending,calc,patterns,positioning}
\definecolor{ink}{HTML}{243745}
\definecolor{blue}{HTML}{275DAB}
\definecolor{green}{HTML}{177C63}
\definecolor{red}{HTML}{BD4935}
\definecolor{orange}{HTML}{AA6C16}
\definecolor{purple}{HTML}{715393}
\tikzset{every node/.style={font=\fontsize{9.3}{12.5}\selectfont,text=ink},
axis/.style={-{Stealth[length=4pt]},line width=.65pt,draw=ink!65},
vec/.style={-{Stealth[length=5pt]},line width=1pt,draw=blue},
short/.style={-{Stealth[length=3pt]},line width=.8pt},
guide/.style={dashed,draw=ink!40,line width=.5pt},
card/.style={rounded corners=2pt,draw=blue!65,fill=blue!4,align=center,text width=6.5cm,inner xsep=9pt,inner ysep=7pt},
title/.style={font=\fontsize{10.5}{14}\selectfont\bfseries,text width=15cm,align=center},
label/.style={font=\fontsize{8.4}{11}\selectfont,align=center},
note/.style={label,text width=6.2cm}}
\begin{document}
\begin{tikzpicture}[x=1cm,y=1cm]
\path[use as bounding box] (-8,-3.5) rectangle(8,3.5);
\node[title] at (0,3.15) {M3-016\quad 介质平面上的带电导体球};
\begin{scope}[shift={(-4,0)}]
\fill[green!5] (-2.5,0) rectangle(2.5,2.2);
\fill[blue!5] (-2.5,-2.2) rectangle(2.5,0);
\draw[ink!60] (-2.5,0)--(2.5,0);
\draw[green,dashed] (0,0)circle(1.9);
\filldraw[draw=ink!65,fill=ink!12] (0,0)circle(.9);
\node at(0,0){$Q$};\draw[guide](0,0)--(140:.9) node[midway,above] {$R$};
\foreach \a in {20,65,115,160,200,245,295,340}{\draw[vec] (\a:1.0)--(\a:1.65);}
\node[label] at(-1.75,1.95){$\varepsilon_{r1},S_1$};\node[label] at(-1.75,-1.95){$\varepsilon_{r2},S_2$};
\node[label] at(2.1,-.3){$r$};\node[label] at(0,-2.6){界面过球心；每个半球面积 $2\pi r^2$};\end{scope}
\node[card,draw=blue!65,fill=blue!4] at (3.9,1.8) {界面上 $E_{1n}=E_{2n}=0$\\$E_{1t}=E_{2t}$};
\node[card,draw=green!65,fill=green!4] at (3.9,0) {$2\pi r^2(D_1+D_2)=Q$};
\node[card,draw=orange!65,fill=orange!4] at (3.9,-1.8) {$\displaystyle\mathbf E=\frac{Q\mathbf e_r}{2\pi\varepsilon_0(\varepsilon_{r1}+\varepsilon_{r2})r^2}$};
\end{tikzpicture}\end{document}